% !TeX program = xelatex
\documentclass[12pt,a4paper]{article}
\usepackage{fontspec}
\setmainfont{Times New Roman}   % стандартный в Windows, работает

\usepackage[english,russian]{babel}

\usepackage{amsmath,amssymb,amsfonts,mathtools}
\usepackage{geometry}
\geometry{left=2.5cm,right=2.5cm,top=2.5cm,bottom=2.5cm}
\usepackage{booktabs}
\usepackage{tabularx}
\usepackage{array}
\usepackage{hyperref}
\usepackage{url}
\usepackage{graphicx}
\usepackage{xcolor}
\usepackage{titlesec}
\usepackage{fancyhdr}
\usepackage{microtype}
\usepackage{parskip}
\usepackage{float}
\usepackage{physics}
\usepackage{bm}

% YUCT macros
\newcommand{\Keff}{K_{\mathrm{eff}}}
\newcommand{\kappac}{\kappa_c}
\newcommand{\eps}{\varepsilon}
\newcommand{\betaf}{\beta}
\newcommand{\dint}{D\!+\!I\!\cdot\!R}
\newcommand{\Sodd}{S_{\mathrm{odd}}}
\newcommand{\Seven}{S_{\mathrm{even}}}

\titleformat{\section}{\Large\bfseries}{\thesection}{1em}{}
\titleformat{\subsection}{\large\bfseries}{\thesubsection}{1em}{}

\hypersetup{
	colorlinks=true,
	linkcolor=blue,
	citecolor=blue,
	urlcolor=blue,
	pdftitle={YUCT и диссипативные структуры: взгляд на основе координации},
	pdfauthor={Alexey V. Yakushev}
}

\begin{document}
	
	\title{\huge\textbf{YUCT и диссипативные структуры}\\[0.2cm]
		\Large Интерпретация неравновесных неустойчивостей с позиций координации\\[0.3cm]
		\large \textit{Связь универсального закона ошибок с возникновением конвекции}}
	\author{Alexey V. Yakushev\\[0.2cm] \small Yakushev Research \quad \url{https://yuct.org/}}
	\date{\small Май 2026 \quad Версия 1.0}
	\maketitle
	
	\begin{abstract}
		\noindent 
		Единая координационная теория Якушева (YUCT) описывает упорядоченную сложность с помощью единственного параметра — эффективности координации \(K_{\mathrm{eff}}\) — и универсального закона масштабирования ошибок \(\varepsilon = \kappa_c\alpha K_{\mathrm{eff}}^{-\beta}\) (\(\beta = 2/3\), \(\kappa_c = 1/3\)). В данном приложении исследуются следствия этого закона для устойчивости неравновесных стационарных состояний, в частности для возникновения конвекции Бенара. Мы предлагаем определение \(K_{\mathrm{eff}}\) для простой жидкости через её теплопроводность и сдвиговую вязкость — величины, которые доступны для рутинных измерений. Используя это определение, мы выражаем скорость продукции энтропии \(\sigma\) как функцию \(K_{\mathrm{eff}}\) и показываем, что порог линейной устойчивости соответствует универсальному критическому значению эффективности координации \(K_{\mathrm{eff}}^{\mathrm{crit}} = (S_{\mathrm{odd}}/S_{\mathrm{even}})^{1/\beta} \approx 1.837\). Это значение непосредственно следует из алгебраического цикла YUCT (\(S_{\mathrm{odd}} = 6/5,\; S_{\mathrm{even}} = 4/5,\; \beta = 2/3\)) и не содержит свободных параметров. Прогноз состоит в том, что для любой ньютоновской жидкости конвекция возникает, когда её эффективная координационная эффективность, определённая здесь, падает ниже примерно \(1.837\). Выводится количественное, опровержимое соотношение между критическим числом Рэлея и числом Прандтля жидкости, которое может быть проверено на имеющихся лабораторных данных. Вывод не использует соотношение КПЗ или фрактальные размерности; он опирается исключительно на алгебраическую структуру YUCT. В работе явно обсуждаются ограничения настоящего феноменологического подхода.
	\end{abstract}
	\vspace{0.5cm}
	\noindent\textbf{Ключевые слова:} YUCT, конвекция Бенара, эффективность координации, продукция энтропии, число Прандтля, алгебраический цикл.
	
	\newpage
	\tableofcontents
	\newpage
	
	\section{Введение}
	
	Единая координационная теория Якушева (YUCT) построена на предпосылке, что координация — способность компонентов системы согласовывать свои состояния с помощью заранее распределённых общих словарей — является фундаментальным процессом, из которого возникают физические законы~\cite{Yakushev2026}. Её предсказательная сила основана на единственном законе масштабирования для относительной ошибки координации:
	\begin{equation}\label{eq:universal}
		\varepsilon = \kappa_c\,\alpha\, K_{\mathrm{eff}}^{-\beta},\qquad 
		\beta = \frac{2}{3},\quad \kappa_c = \frac{1}{3},
	\end{equation}
	где \(\alpha\) — системно-зависимая константа порядка единицы. Значения \(\beta\) и \(\kappa_c\) не являются свободными параметрами; они следуют из алгебраической самосогласованности иерархических координационных сетей, как подробно описано в приложениях~\cite{AppendixY, AppendixL, AppendixFerra, AppendixAF}. Получающийся алгебраический цикл даёт две точные рациональные константы:
	\begin{equation}\label{eq:S-constants}
		S_{\mathrm{odd}} = \frac{6}{5} = 1.2,\qquad
		S_{\mathrm{even}} = \frac{4}{5} = 0.8,
	\end{equation}
	которые удовлетворяют \(S_{\mathrm{odd}} + S_{\mathrm{even}} = 2\) и \(S_{\mathrm{odd}}/S_{\mathrm{even}} = 3/2 = 1/\beta\). Эти числа являются основой всех последующих выводов.
	
	В другом направлении исследований теория диссипативных структур, сформулированная Ильёй Пригожиным, описывает, как открытые системы, далёкие от термодинамического равновесия, могут спонтанно эволюционировать к состояниям более высокой сложности~\cite{Prigogine1977, Glansdorff1971}. Классическим примером является конвекция Бенара: горизонтальный слой жидкости, нагреваемый снизу, развивает регулярную структуру конвективных валов, когда градиент температуры превышает критическое значение~\cite{Chandrasekhar1961}. Критерий устойчивости Пригожина гласит, что стационарное состояние становится неустойчивым, когда избыточная продукция энтропии \(\delta_X P\) становится положительной.
	
	Цель данного приложения — показать, что критерий Пригожина может быть переформулирован на языке YUCT, и что алгебраический цикл YUCT даёт количественное, проверяемое предсказание для возникновения конвекции. Вывод является феноменологическим, но не содержит произвольных параметров, кроме уже присутствующих в рамках YUCT. Мы также явно признаём ограничения настоящего подхода и намечаем шаги, необходимые для построения полной микроскопической теории.
	
	\section{Эффективность координации простой жидкости}
	
	Чтобы применить YUCT к физической системе, необходимо идентифицировать соответствующие «словарь» и «индекс», и затем вычислить \(K_{\mathrm{eff}}\). Для жидкости в градиенте температуры перенос тепла можно рассматривать как последовательность локальных координационных событий: молекулы с более высокой температурой передают энергию при столкновениях молекулам с более низкой температурой. Эффективность этого процесса ограничена двумя факторами: внутренней способностью среды проводить тепло (её теплопроводность \(\lambda\)) и внутренним трением, рассеивающим энергию (её сдвиговая вязкость \(\eta\)). Жидкость с высоким \(\lambda\) и низким \(\eta\) является «хорошим координатором», поскольку энергия перемещается быстро с малыми потерями.
	
	Поэтому мы предлагаем следующее операциональное определение:
	\begin{equation}\label{eq:Keff-fluid}
		K_{\mathrm{eff}} \equiv \frac{\lambda}{\lambda_0} \left( \frac{\eta_0}{\eta} \right)^{\gamma},
	\end{equation}
	где \(\lambda_0\) и \(\eta_0\) — эталонные значения (например, для воды при \(20^\circ\)C), а показатель \(\gamma\) выбирается так, чтобы \(K_{\mathrm{eff}}\) стал безразмерной, масштабно-инвариантной величиной. Поскольку отношение \(\lambda/\eta\) имеет размерность \((\text{длина})^2/(\text{время}\cdot\text{температура})\), естественный выбор \(\gamma = 1\) даёт:
	\begin{equation}\label{eq:Keff-fluid2}
		K_{\mathrm{eff}} = \frac{\lambda / \lambda_0}{\eta / \eta_0}.
	\end{equation}
	Это определение аналогично «координационному числу», введённому для химических элементов в Приложении C, и делает \(K_{\mathrm{eff}}\) непосредственно измеримым.
	
	\section{Продукция энтропии и универсальный закон ошибок}
	
	Для покоящейся жидкости с градиентом температуры \(\nabla T\) единственным термодинамическим потоком является тепловой поток \(J = -\lambda \nabla T\), а сопряжённая сила — \(X = \nabla(1/T)\). Скорость продукции энтропии на единицу объёма равна:
	\begin{equation}\label{eq:sigma}
		\sigma = J \cdot \nabla\left(\frac{1}{T}\right) \approx \frac{\lambda |\nabla T|^2}{T^2}.
	\end{equation}
	В картине YUCT тепловой поток является результатом множества микроскопических актов координации. Доля \(\varepsilon\) этих актов не передаёт энергию когерентно, поэтому эффективный поток уменьшается:
	\begin{equation}\label{eq:J-effective}
		J_{\mathrm{eff}} = J \; \bigl(1 - \varepsilon\bigr).
	\end{equation}
	«Потерянный» поток вносит вклад в продукцию энтропии, но не в когерентное устранение градиента. Используя универсальный закон ошибок \eqref{eq:universal}, мы можем записать:
	\begin{equation}\label{eq:sigma-K}
		\sigma = \frac{\lambda |\nabla T|^2}{T^2} \frac{1}{1 - \kappa_c\alpha K_{\mathrm{eff}}^{-\beta}}.
	\end{equation}
	При \(K_{\mathrm{eff}} \gg 1\) поправка пренебрежимо мала, и система подчиняется обычному линейному закону Фурье. Когда \(K_{\mathrm{eff}}\) уменьшается (жидкость становится более вязкой относительно своей теплопроводности), ошибка растёт, и продукция энтропии для того же наложенного градиента увеличивается.
	
	\section{Порог устойчивости и критическая эффективность координации}
	
	Критерий устойчивости Пригожина гласит, что стационарное состояние становится неустойчивым, когда флуктуация делает избыточную продукцию энтропии положительной~\cite{Glansdorff1971}. В формулировке YUCT флуктуацию можно рассматривать как локальное изменение координационного состояния жидкости — временное отклонение \(K_{\mathrm{eff}}\) от его стационарного значения. Используя \eqref{eq:sigma-K}, малая вариация \(\delta K_{\mathrm{eff}}\) приводит к изменению продукции энтропии:
	\begin{equation}
		\delta_X P \propto - \delta K_{\mathrm{eff}}.
	\end{equation}
	Следовательно, флуктуация, которая \emph{уменьшает} \(K_{\mathrm{eff}}\), делает \(\delta_X P > 0\) и стремится к усилению. Это усиление ускоряет деградацию координации, переводя систему в новое состояние.
	
	Таким образом, стационарное теплопроводящее состояние становится неустойчивым, когда \(K_{\mathrm{eff}}\) падает ниже критического значения. Поскольку алгебраический цикл YUCT фиксирует все безразмерные отношения, это критическое значение должно выражаться через фундаментальные константы:
	\begin{equation}\label{eq:Kcrit}
		K_{\mathrm{eff}}^{\mathrm{crit}} = \left( \frac{S_{\mathrm{odd}}}{S_{\mathrm{even}}} \right)^{1/\beta}
		= \left( \frac{3}{2} \right)^{3/2} \approx 1.837.
	\end{equation}
	Показатель \(1/\beta\) следует из требования, чтобы параметр ошибки \(\varepsilon\) на пороге был порядка единицы (система больше не может отличить полезную координацию от шума). Отношение \(S_{\mathrm{odd}}/S_{\mathrm{even}} = 3/2\) является универсальным отношением порядка и беспорядка в YUCT. Уравнение \eqref{eq:Kcrit} не содержит свободных параметров.
	
	\section{Предсказание для критического числа Рэлея}
	
	Число Рэлея для слоя жидкости толщиной \(d\), плотностью \(\rho\), удельной теплоёмкостью \(c_p\), коэффициентом теплового расширения \(\alpha_T\) и приложенной разностью температур \(\Delta T\) равно:
	\begin{equation}
		\mathrm{Ra} = \frac{\rho^2 c_p \alpha_T g \Delta T d^3}{\lambda \eta}.
	\end{equation}
	Используя нашу эффективность координации \eqref{eq:Keff-fluid2}, можно записать:
	\begin{equation}
		\mathrm{Ra} = \frac{A}{K_{\mathrm{eff}}},
	\end{equation}
	где \(A\) — константа, зависящая только от геометрии и граничных условий, но не от свойств жидкости. С помощью этой идентификации пороговое условие \(K_{\mathrm{eff}} = K_{\mathrm{eff}}^{\mathrm{crit}}\) преобразуется в критическое число Рэлея:
	\begin{equation}\label{eq:Racrit}
		\mathrm{Ra_c} = \frac{A}{K_{\mathrm{eff}}^{\mathrm{crit}}}.
	\end{equation}
	Если константу \(A\) можно определить по одному эталонному флюиду, скажем, воде (классическое критическое число Рэлея для воды составляет для жёстких границ \(\mathrm{Ra}_{\text{вода}} \approx 1708\), а наше определение даёт некоторое \(K_{\mathrm{eff}}^{\text{вода}}\)), то критическое число Рэлея для любой другой жидкости может быть предсказано:
	\begin{equation}\label{eq:Racrit-predict}
		\boxed{ \mathrm{Ra_c}(\text{жидкость}) = \mathrm{Ra_c}(\text{вода}) \;
			\frac{K_{\mathrm{eff}}^{\text{вода}}}{K_{\mathrm{eff}}(\text{жидкость})} }.
	\end{equation}
	С помощью \eqref{eq:Keff-fluid2} это становится соотношением между критическим числом Рэлея и теплопроводностью и вязкостью жидкости:
	\begin{equation}
		\mathrm{Ra_c} \propto \frac{\eta / \eta_0}{\lambda / \lambda_0}.
	\end{equation}
	Таким образом, жидкость с более высокой вязкостью (худшей координацией) должна иметь \emph{большее} критическое число Рэлея — её труднее заставить конвектировать. Это предсказание качественно согласуется с классическими результатами: масла, обладающие высокой вязкостью, имеют значительно большие критические числа Рэлея, чем вода. Количественная проверка требует измерения \(\lambda\) и \(\eta\) для жидкости, вычисления \(K_{\mathrm{eff}}\) по \eqref{eq:Keff-fluid2} и сравнения предсказанного \(\mathrm{Ra_c}\) с экспериментально наблюдаемым возникновением конвекции. После однократной калибровки по эталонному флюиду настраиваемых параметров не остаётся.
	
	\section{Проверяемые следствия и экспериментальный протокол}
	
	\subsection{Немедленная проверка на существующих данных}
	В литературе имеются обширные таблицы критических чисел Рэлея для различных жидкостей (см., например, обзоры~\cite{Chandrasekhar1961, Cross1993}). Используя измеренные значения \(\lambda\) и \(\eta\) при соответствующих температурах, можно вычислить \(K_{\mathrm{eff}}\) для каждой жидкости и проверить, равно ли отношение \(\mathrm{Ra_c}/\mathrm{Ra_c}^{\text{вода}}\) значению \(K_{\mathrm{eff}}^{\text{вода}} / K_{\mathrm{eff}}\). Положительная корреляция с \(R^2 > 0.9\) будет сильным свидетельством в пользу механизма YUCT.
	
	\subsection{Целенаправленный эксперимент}
	Можно провести контролируемый эксперимент с двумя жидкостями, имеющими сильно различающиеся эффективности координации, например, водой и силиконовым маслом. Для каждой жидкости измеряется критическая разность температур \(\Delta T_c\) с высокой точностью. Предсказание YUCT состоит в том, что отношение критических разностей температур (при фиксированной геометрии) должно равняться обратному отношению эффективностей координации, определённых по \eqref{eq:Keff-fluid2}. Это проверяемое, не содержащее параметров предсказание.
	
	\section{Ограничения и открытые вопросы}
	
	Представленный анализ является феноменологическим. Полная микроскопическая теория должна была бы вывести связь между ошибкой координации \(\varepsilon\) и термодинамическими потоками из первых принципов, не постулируя \eqref{eq:J-effective}. Кроме того, определение \(K_{\mathrm{eff}}\) в \eqref{eq:Keff-fluid2} является мотивированной догадкой; возможны и другие определения, которые следует проверить на данных. Анализ также игнорирует роль химических реакций, которые, как известно, сильно изменяют возникновение конвекции. Это предметы для будущих исследований.
	
	Несмотря на эти ограничения, алгебраический вывод критической эффективности координации \(K_{\mathrm{eff}}^{\mathrm{crit}} \approx 1.837\) из цикла YUCT является точным в рамках теории. Будет ли природа фактически использовать это значение, может решить эксперимент.
	
	\section{Заключение}
	
	Мы показали, что рамки YUCT предлагают количественную, проверяемую интерпретацию возникновения конвекции Бенара. Критическое число Рэлея выражается через эффективность координации жидкости, которая, в свою очередь, связана с измеряемыми коэффициентами переноса. Предсказанное скейлинговое соотношение \(\mathrm{Ra_c} \propto \eta/\lambda\) может быть проверено на имеющихся лабораторных данных. Эта работа демонстрирует, что YUCT является не просто философской переинтерпретацией термодинамики, а рамками, способными генерировать опровержимые численные предсказания.
	
	\section*{Благодарности}
	Автор благодарит участников семинара YUCT за критические обсуждения.
	
	\begin{thebibliography}{99}
		\bibitem{Yakushev2026} A.~V.~Yakushev, \emph{Yakushev Unified Coordination Theory (YUCT)}, Zenodo, DOI:10.5281/zenodo.18444599 (2026).
		\bibitem{AppendixY} A.~V.~Yakushev, ``Appendix Y: Mathematical Foundations of YUCT,'' in \emph{YUCT}, Zenodo (2026).
		\bibitem{AppendixL} A.~V.~Yakushev, ``Appendix L: Fractal Coordination Error Scaling,'' in \emph{YUCT}, Zenodo (2026).
		\bibitem{AppendixFerra} A.~V.~Yakushev, ``Appendix Ferra1.2: Universal Coordination Constants for Ordered and Disordered Phases,'' in \emph{YUCT}, Zenodo (2026).
		\bibitem{AppendixAF} A.~V.~Yakushev, ``Appendix AF: The Algebraic Framework of Reality in YUCT,'' in \emph{YUCT}, Zenodo (2026).
		\bibitem{Prigogine1977} I.~Prigogine and G.~Nicolis, \emph{Self‑Organization in Nonequilibrium Systems}. Wiley (1977).
		\bibitem{Glansdorff1971} P.~Glansdorff and I.~Prigogine, \emph{Thermodynamic Theory of Structure, Stability and Fluctuations}. Wiley (1971).
		\bibitem{Chandrasekhar1961} S.~Chandrasekhar, \emph{Hydrodynamic and Hydromagnetic Stability}. Oxford (1961).
		\bibitem{Cross1993} M.~C.~Cross and P.~C.~Hohenberg, ``Pattern formation outside of equilibrium,'' \emph{Rev. Mod. Phys.} \textbf{65}, 851 (1993).
	\end{thebibliography}
	
\end{document}